% glossary.tex -- Module 5 glossary (A-Z). GENERATED by
% not this file.
\unnumberedchapter[\faBook]{Glossary}

Every technical word in this book, from A to Z. The meaning is in plain English. The
last column gives the usual Chinese term used in physics textbooks.

\begin{glossarytable}
  % ---- A
  \gterm{Acceleration due to gravity ($g$)}{the acceleration of a falling object near the Earth: \qty{9.8}{m.s^{-2}} downwards}{重力加速度}
  \gterm{Accuracy}{how close a result is to the true or accepted value}{准确度}
  \gterm{Air resistance (drag)}{the push of the air on a moving object, opposite to its velocity}{空气阻力}
  \gterm{Altitude ($h$)}{the height above a planet's surface}{离地高度}
  \gterm{Angular displacement ($\Delta\theta$)}{the angle an object turns through, measured at the centre}{角位移}
  \gterm{Angular velocity ($\omega$)}{the angle turned each second, in \si{rad.s^{-1}}}{角速度}
  \gterm{Anticlockwise}{turning the opposite way to the hands of a clock}{逆时针}
  \gterm{Aphelion}{the farthest point of an orbit from the Sun}{远日点}
  \gterm{Apogee}{the farthest point of an orbit around the Earth}{远地点}
  \gterm{Apparent weight}{the push of the seat or floor on you; it is what you feel as your weight}{视重}
  \gterm{Apparent weightlessness}{the feeling of having no weight in free fall, because the normal force is zero}{完全失重}
  \gterm{Arc length ($s$)}{the distance along the edge of a circle: $s = r\theta$, with $\theta$ in radians}{弧长}
  \gterm{Assumption}{something we accept as true to keep the maths simple}{假设}
  \gterm{Astronomical unit (AU)}{the average distance from the Earth to the Sun, about \qty{1.50e11}{m}}{天文单位}
  \gterm{Atmospheric drag}{the push of thin air against a satellite; it turns mechanical energy into heat}{大气阻力}
  % ---- B
  \gterm{Balanced torques}{clockwise and anticlockwise torques are equal, so the rotation does not change}{力矩平衡}
  \gterm{Bank angle}{the angle between a banked road and the horizontal}{倾角}
  \gterm{Banked track}{a curved road that is tilted, higher on the outside of the bend}{倾斜弯道（外高内低）}
  \gterm{Bound}{trapped by gravity, so the object cannot get away; its total energy is negative}{被引力束缚的}
  \gterm{Burn (engine burn)}{a short time when a rocket engine fires to change a spacecraft's speed}{点火（发动机点火）}
  % ---- C
  \gterm{Centre of mass}{the point where all of an object's mass (and its weight) seems to act}{质心}
  \gterm{Centrifugal force}{a made-up ``outward force''; do not use it: the outward feeling is inertia}{离心力}
  \gterm{Centripetal acceleration ($a_\mathrm{c}$)}{the acceleration towards the centre of a circle: $v^2/r$}{向心加速度}
  \gterm{Centripetal force ($F_\mathrm{c}$)}{the net force towards the centre that keeps an object moving in a circle}{向心力}
  \gterm{Circumference}{the distance once around a circle, $2\pi r$}{周长}
  \gterm{Clockwise}{turning the same way as the hands of a clock}{顺时针}
  \gterm{Coefficient of friction ($\mu$)}{a number with no unit that says how grippy two surfaces are}{摩擦因数}
  \gterm{Complementary angles}{two angles that add up to \ang{90}, such as \ang{30} and \ang{60}}{互余角（余角）}
  \gterm{Component}{the part of a vector (such as a force) that points in one chosen direction}{分量（分力）}
  \gterm{Conical pendulum}{a mass on a string moving in a horizontal circle; the string sweeps out a cone}{圆锥摆}
  \gterm{Conservation of energy}{energy is never made or destroyed; it only changes form, so the total stays the same}{能量守恒}
  \gterm{Conservation of mechanical energy}{if only gravity does work, $K + U$ stays the same}{机械能守恒}
  \gterm{Constant acceleration}{an acceleration that does not change, so the velocity changes by the same amount every second}{恒定加速度（匀变速）}
  \gterm{Controlled variable}{a quantity you keep the same, so the test is fair}{控制变量}
  \gterm{Crank}{the arm from a bicycle pedal to the axle}{曲柄}
  \gterm{Crest}{the top of a hill in a track; the centre of the circle is below it}{坡顶}
  % ---- D
  \gterm{Dense}{heavy for its size}{密度大的}
  \gterm{Dependent variable}{the quantity you measure to see the effect}{因变量}
  \gterm{Design speed}{the one speed at which a banked curve needs no friction}{设计速度（规定速度）}
  \gterm{Dip}{the low point of a valley in a track; the centre of the circle is above it}{谷底}
  \gterm{Displacement ($s$)}{the change in position, with a direction}{位移}
  % ---- E
  \gterm{Eccentricity ($e$)}{how stretched an ellipse is: 0 is a circle, close to 1 is very flat}{离心率（偏心率）}
  \gterm{Ellipse}{a stretched circle (an oval)}{椭圆}
  \gterm{Elliptical orbit}{an oval orbit with the planet or the Sun at one focus}{椭圆轨道}
  \gterm{Empirical law}{a rule found from data that describes motion but does not explain it}{经验定律}
  \gterm{Escape velocity ($v_\mathrm{esc}$)}{the smallest launch speed that lets a projectile get away from a planet for ever}{逃逸速度}
  % ---- F
  \gterm{Field lines}{arrows that show which way a mass would be pulled in a field}{场线}
  \gterm{Final velocity}{the velocity just before landing: a speed and a direction}{末速度}
  \gterm{Focus (plural: foci)}{one of two fixed points inside an ellipse; the Sun is at one focus of a planet's orbit}{焦点}
  \gterm{Free fall}{motion when gravity is the only force acting}{自由下落（只受引力）}
  \gterm{Free-body diagram}{a sketch of one object with an arrow for each real force on it}{受力分析图}
  \gterm{Frequency ($f$)}{the number of full turns (or cycles) each second, in hertz (\si{Hz})}{频率}
  % ---- G
  \gterm{Geostationary orbit}{a circular orbit above the equator with a one-day period, so the satellite stays above one point}{地球静止轨道}
  \gterm{Geosynchronous orbit}{any orbit with a period of one day; it may be tilted}{地球同步轨道}
  \gterm{Gradient}{rise divided by run for a straight line on a graph}{斜率}
  \gterm{Gravitational field}{the region around a mass where another mass feels a gravitational force}{引力场}
  \gterm{Gravitational field strength ($g$)}{the gravitational force on each kilogram at a point, in \si{N.kg^{-1}} (\si{m.s^{-2}})}{引力场强度}
  \gterm{Gravitational force}{the pull between any two objects that have mass}{引力（万有引力）}
  \gterm{Gravitational potential energy ($U$)}{energy stored because of where an object is in a gravitational field; near the ground it changes by $mg\Delta h$; in space $U = -GMm/r$}{引力势能（地面附近：重力势能）}
  % ---- H
  \gterm{Hohmann transfer orbit}{a fuel-saving half-ellipse between two circular orbits, using two forward engine burns}{霍曼转移轨道}
  \gterm{Hypotenuse}{the longest side of a right-angled triangle}{斜边}
  % ---- I
  \gterm{Independent (motions)}{not affecting each other; the horizontal and vertical motions share only the time}{相互独立}
  \gterm{Independent variable}{the quantity you change on purpose in an experiment}{自变量}
  \gterm{Inertia}{the tendency of an object to keep moving in a straight line at a steady speed}{惯性}
  \gterm{Infinity}{a place so far away that the planet's gravity there is too small to matter}{无穷远处}
  \gterm{Inverse-square law}{a rule where doubling the distance makes the effect four times smaller}{平方反比定律}
  % ---- K
  \gterm{Kepler's first law (law of ellipses)}{each planet moves in an ellipse with the Sun at one focus}{开普勒第一定律（轨道定律）}
  \gterm{Kepler's second law (law of equal areas)}{the line from the Sun to a planet sweeps out equal areas in equal times}{开普勒第二定律（面积定律）}
  \gterm{Kepler's third law (law of periods)}{$T^2$ is proportional to $r^3$ for everything that orbits the same body}{开普勒第三定律（周期定律）}
  \gterm{Kinetic energy ($K$)}{energy of motion: $K = \tfrac{1}{2}mv^2$}{动能}
  % ---- L
  \gterm{Launch angle ($\theta$)}{the angle of the launch velocity above the horizontal}{抛射角}
  \gterm{Launch height ($h$)}{how far the launch point is above the landing point}{抛出点高度}
  \gterm{Level ground}{landing at the same height as the launch point}{水平地面（落点与抛出点等高）}
  \gterm{Lever}{a bar that turns about a pivot, such as a spanner or a door}{杠杆}
  \gterm{Light gate}{a light beam and timer that measures how long an object blocks the beam}{光电门}
  \gterm{Line of action}{the line along which a force acts, drawn longer in both directions}{力的作用线}
  \gterm{Line of best fit}{the straight line that passes as close as possible to all the data points}{最佳拟合直线}
  \gterm{Low Earth orbit (LEO)}{an orbit about \qtyrange{200}{2000}{km} above the Earth, with a period of about 90 minutes to 2 hours}{低地球轨道（近地轨道）}
  % ---- M
  \gterm{Mass ($m$)}{the amount of matter in an object, in kg; the same everywhere}{质量}
  \gterm{Maximum height ($h_\mathrm{max}$)}{the greatest height above the launch point}{最大高度}
  \gterm{Model}{a simple description of real life that we can calculate with}{模型}
  % ---- N
  \gterm{Net force}{the total of all the forces on an object, added with their directions}{合力}
  \gterm{Non-uniform circular motion}{moving in a circle with a changing speed}{变速圆周运动}
  \gterm{Normal force ($F_\mathrm{N}$)}{the push of a surface on an object, at \ang{90} to the surface}{支持力（法向力）}
  % ---- O
  \gterm{Orbit}{the path of a satellite around a planet, or of a planet around a star}{轨道}
  \gterm{Orbital decay}{the slow fall to lower orbits caused by drag from the thin upper atmosphere}{轨道衰减}
  \gterm{Orbital period ($T$)}{the time for one full orbit}{轨道周期}
  \gterm{Orbital radius ($r$)}{the distance from the centre of the planet to the satellite}{轨道半径}
  \gterm{Orbital speed ($v$)}{the speed a satellite needs to stay in a circular orbit}{环绕速度（轨道速度）}
  \gterm{Origin}{the point we measure every position from, $(0, 0)$}{原点}
  % ---- P
  \gterm{Parabola}{a curve that is symmetrical about its highest point; the path of a projectile}{抛物线}
  \gterm{Parking orbit}{a low circular orbit where a spacecraft waits before its next engine burn}{停泊轨道}
  \gterm{Perigee}{the closest point of an orbit around the Earth}{近地点}
  \gterm{Perihelion}{the closest point of an orbit to the Sun}{近日点}
  \gterm{Period ($T$)}{the time for one full turn}{周期}
  \gterm{Perpendicular distance ($r_\perp$)}{the shortest distance from the pivot to the line of action of a force (the lever arm)}{力臂}
  \gterm{Pivot}{the fixed point or line that an object turns about, or that a string hangs from}{支点（转轴；悬点）}
  \gterm{Plumb line}{a weighted string that hangs straight down}{铅垂线}
  \gterm{Point mass}{an object treated as if all its mass is at one point}{质点}
  \gterm{Prediction}{a result you calculate from a model before you measure it}{预测}
  \gterm{Primary data}{data you measure yourself}{一手数据（原始数据）}
  \gterm{Projectile}{an object that is thrown or launched and then moves with no engine; in the model, gravity is the only force on it}{抛体}
  \gterm{Proportional}{two quantities are proportional when doubling one doubles the other}{成正比}
  % ---- Q
  \gterm{Quadratic equation}{an equation with a $t^2$ term, a $t$ term and a number ($At^2 + Bt + C = 0$)}{一元二次方程}
  % ---- R
  \gterm{Radial field}{a field whose lines all point straight at the centre of the planet, like the spokes of a wheel}{径向场}
  \gterm{Radian (rad)}{a unit for angles: one full turn is $2\pi$ rad, the same as \ang{360}}{弧度}
  \gterm{Random error}{an error that changes from trial to trial, sometimes too big and sometimes too small}{偶然误差（随机误差）}
  \gterm{Range ($R$)}{the horizontal distance from the launch point to the landing point}{射程（水平射程）}
  \gterm{Reliability}{whether repeated trials give the same result}{可靠性}
  \gterm{Revolutions per minute (rpm)}{the number of full turns in one minute}{转每分（r/min）}
  \gterm{Root (of an equation)}{a solution of an equation}{根（方程的解）}
  \gterm{Rotor (ride)}{a spinning drum ride whose riders stay on the wall when the floor drops away}{旋转圆筒（转筒）}
  % ---- S
  \gterm{Satellite}{an object that orbits a bigger body; the Moon is a natural satellite}{卫星}
  \gterm{Scalar}{a quantity with a size only, such as speed, time, mass or energy}{标量}
  \gterm{Scientific notation}{a number from 1 to 10, times a power of ten, such as \num{2.9e2}}{科学记数法}
  \gterm{Secondary data}{measurements made by other people, such as space agencies}{二手数据}
  \gterm{Semi-major axis ($a$)}{half of the longest diameter of an ellipse}{半长轴}
  \gterm{Significant figures}{the digits in a number that carry meaning}{有效数字}
  \gterm{Speed}{the size of the velocity, with no direction}{速率}
  \gterm{Static friction ($F_\mathrm{f}$)}{friction that stops surfaces sliding; it grows as needed, up to a maximum of $\mu F_\mathrm{N}$}{静摩擦力}
  \gterm{Symmetrical}{the same on both sides; the way down mirrors the way up}{对称的}
  \gterm{Systematic error}{an error in the same direction every time; repeating does not remove it}{系统误差}
  % ---- T
  \gterm{Tangent}{a line that touches a circle at one point, at \ang{90} to the radius}{切线}
  \gterm{Tangential speed ($v$)}{the speed of a point as it moves round its circle: $v = \omega r$}{线速度}
  \gterm{Tension ($F_\mathrm{T}$)}{the pull of a string, rope or chain, along its length}{拉力（张力）}
  \gterm{Time of flight ($t_\mathrm{flight}$)}{the time from launch until landing}{飞行时间}
  \gterm{Time step}{how far forward in time a computer model jumps in each calculation}{时间步长}
  \gterm{Torque ($\tau$)}{the turning effect of a force, in \si{N.m}}{力矩}
  \gterm{Total mechanical energy ($E$)}{kinetic energy plus potential energy: $E = K + U$}{机械能}
  \gterm{Trajectory}{the path that a moving object follows}{轨迹}
  \gterm{Transfer orbit}{the elliptical path that joins two circular orbits}{转移轨道}
  % ---- U
  \gterm{Uniform circular motion}{moving around a circle at a constant speed}{匀速圆周运动}
  \gterm{Uniform field}{a field with the same size and direction everywhere}{匀强场}
  \gterm{Uniform velocity}{a velocity that does not change: the same speed in the same direction; the acceleration is zero}{匀速（速度恒定）}
  \gterm{Universal gravitation}{Newton's idea that one law of gravity works everywhere, on Earth and in space}{万有引力定律}
  \gterm{Universal gravitational constant ($G$)}{the fixed number that sets how strong gravity is, \qty{6.67e-11}{N.m^2.kg^{-2}}}{引力常量}
  % ---- V
  \gterm{Validate}{to test whether predictions match real measurements}{验证}
  \gterm{Validity}{whether an experiment really tests what it claims to test}{有效性}
  \gterm{Vector}{a quantity with a size and a direction, drawn as an arrow}{矢量}
  \gterm{Vertical circle}{a circle in an upright plane, like the path of a swing; the speed changes}{竖直平面内的圆周}
  % ---- W
  \gterm{Weight}{the pull of gravity on an object: $mg$, straight down, in N}{重力}
  \gterm{Work ($W$)}{energy transferred by a force when the object moves along the direction of the force (or part of it)}{功}
  \gterm{Work against gravity}{the energy an engine gives an object to move it away from a planet; it is stored as $U$}{克服引力做的功}
  \gterm{Work--energy theorem}{the work done by the net force equals the change in kinetic energy}{动能定理}
\end{glossarytable}
